% !TeX program = xelatex
\documentclass[12pt]{nextbook_zh}

\UseTemplateSet{
  globalfonts = {libertinus,japanese},
  features = {hyperlinks, tables}
}

\newcommand{\ReaderVocabBaselineskip}{15pt}
\usepackage{paracol}
\usepackage{xcolor}
\usepackage{eso-pic}
\usepackage{needspace}

\vbadness=10000
\hbadness=10000
\hfuzz=3pt
\emergencystretch=1.5em

\makeatletter
\AtBeginDocument{%
  \@ifpackageloaded{hyperref}{%
    \pdfstringdefDisableCommands{%
      \def\ruby#1#2{#1}%
      \def\ReaderVocabText#1{#1}%
      \def\ReaderVocabTerm#1{#1}%
    }%
  }{}%
}
\makeatother

\providecommand{\ReaderColumnSep}{6mm}
\providecommand{\ReaderColumnRatio}{0.75}
\providecommand{\ReaderVocabBaselineskip}{20pt}
\providecommand{\ReaderVocabText}[1]{\textbf{#1}}
\providecommand{\ReaderVocabTerm}[1]{\textbf{#1}}
\providecommand{\ReaderVocabSeparator}{:}
\providecommand{\ReaderVocabItem}[2]{\noindent\ReaderVocabTerm{#1}\ReaderVocabSeparator\ #2\par}

\setlength{\columnsep}{\ReaderColumnSep}
\setlength{\columnseprule}{0pt}
\columnratio{\ReaderColumnRatio}
\swapcolumninevenpages

\newcounter{readersection}
\newwrite\vocabfile
\newif\iffirstreaderlabelwritten
\newlength{\readercolumnswidth}
\newlength{\readertitlegutter}
\setlength{\readercolumnswidth}{\dimexpr\textwidth-\columnsep\relax}
\setlength{\readertitlegutter}{.75\columnsep}

\makeatletter
\newcommand{\readerseparatorline}{%
  \if@mainmatter
    \AtTextLowerLeft{%
      \ifodd\value{page}%
        \put(\LenToUnit{\dimexpr .75\readercolumnswidth+.5\columnsep\relax},0){%
          \color{black!35}\rule{0.3pt}{\textheight}%
        }%
      \else
        \put(\LenToUnit{\dimexpr .25\readercolumnswidth+.5\columnsep\relax},0){%
          \color{black!35}\rule{0.3pt}{\textheight}%
        }%
      \fi
    }%
  \fi
}
\makeatother

\AddToShipoutPictureBG{\readerseparatorline}

\newcommand{\currentvocabfile}{\jobname-vocab-\thereadersection.voc}
\newcommand{\beginvocabwrite}{\immediate\openout\vocabfile=\currentvocabfile}
\newcommand{\finishvocabwrite}{\immediate\closeout\vocabfile}
\newcommand{\vocabitem}[2]{\ReaderVocabItem{#1}{#2}}
\newcommand{\printvocablist}{%
  \footnotesize
  \setlength{\parindent}{0pt}%
  \setlength{\parskip}{0pt}%
  \setlength{\baselineskip}{\ReaderVocabBaselineskip}%
  \raggedright
  \input{\currentvocabfile}%
}

\newcommand{\readersectiontitle}[1]{%
  \par
  \begingroup
  \ifodd\value{page}%
    \leftskip=\readertitlegutter%
    \rightskip=0.25\textwidth plus 1fil%
  \else
    \leftskip=\dimexpr .25\readercolumnswidth+.5\columnsep+\readertitlegutter\relax%
    \rightskip=0pt plus 1fil%
  \fi
  \section{#1}%
  \par
  \endgroup
}

\NewDocumentCommand{\voc}{omm}{%
  \ReaderVocabText{#2}%
  \IfNoValueTF{#1}{%
    \immediate\write\vocabfile{\string\vocabitem{#2}{#3}}%
  }{%
    \immediate\write\vocabfile{\string\vocabitem{#1}{#3}}%
  }%
}

\newenvironment{readersection}[1]{%
  \stepcounter{readersection}%
  \beginvocabwrite
  \Needspace{10\baselineskip}%
  \iffirstreaderlabelwritten\else
    \label{reader:firstpage}%
    \global\firstreaderlabelwrittentrue
  \fi
  \readersectiontitle{#1}%
  \begin{paracol}{2}
}{%
  \finishvocabwrite
  \switchcolumn
  \printvocablist
  \end{paracol}
}

\usepackage{ruby}
\renewcommand{\rubysep}{-0.6em}
\let\ReaderOriginalRuby\ruby
\renewcommand{\ruby}[2]{\allowbreak\ReaderOriginalRuby{#1}{#2}\allowbreak}

\title{日本語基礎読本}
\author{Kelvin Yang}
\date{\today}

\begin{document}
\frontmatter
\maketitle
\chapter*{総綱}
{\small
\begin{longtable}{r p{0.52\linewidth} p{0.28\linewidth}}
\toprule
Lesson & Grammar Focus & McGloin Sections \\
\midrule
\endfirsthead

\toprule
Lesson & Grammar Focus & McGloin Sections \\
\midrule
\endhead

\bottomrule
\endfoot

1 & Nominal predicate sentences; X は Y です/だ; minimal sentence structure & 5.1.1; 5.2.1 \\
2 & Adjectival sentences; i-adjectives and na-adjectives as predicates & 5.2.2; 13.1; 13.2 \\
3 & Existential sentences; 場所に N がある/いる and N は場所にある/いる & 5.2.3; 12.3; 10.2.1 \\
4 & Basic verbal sentences; intransitive, transitive, and ditransitive frames; が, を, に & 5.2.4; 10.1.1--10.1.3 \\
5 & Topic-comment sentences; introductory は; the 象は鼻が長い pattern & 5.2.5; 11.1--11.2 \\
6 & の and noun modification; N の N; possession, apposition, and attribution & 7.2.1; 10.1.4 \\
7 & Basic demonstratives; こそあど system; これ, それ, あれ, どれ; ここ, そこ, あそこ & 9.1.1; 9.2 \\
8 & Pronouns; 私, 僕, 彼, 彼女, 自分, お互い; pronoun modification & 8.1; 8.2; 8.3 \\
9 & Interrogatives; 誰, 何, どこ, いつ, どう, どんな & 4.3 \\
10 & Indefinites; 何か, 誰か, 何も, 誰も, 何でも, 誰でも & 4.4 \\
11 & Affirmative and negative forms; nonpast and past; basic polite style & 5.1.2; 5.1.3 \\
12 & Plain and polite styles; spoken and written language; style mixing & 6.1; 6.2 \\
13 & Gendered and rough/gentle styles; ellipsis in conversation; introductory sentence-final particles & 6.3; 1.6; 10.4 \\
14 & Overview of verb conjugation; godan, ichidan, する, 来る; masu-form, dictionary form, te-form, ta-form & 12.1 \\
15 & Verbal negation; 行かない, 食べない, しない, 来ない & 15.1 \\
16 & Inflection of adjectives and nominal predicates; past, negative, and adnominal forms & 13.1.1--13.1.2; 13.2.1--13.2.2 \\
17 & Te-form connection; adjectival te-forms; simple coordination and sequential connection & 22.1; 13.3 \\
18 & Adverbs; frequency, degree, sentential adverbs, and fixed predicate pairings & 14.1; 14.2; 14.3 \\
19 & Numerals and counters; basic quantity expressions & 16.1; 16.2 \\
20 & Quantifier position and frequency expressions; 三冊の本 and 本を三冊 & 16.3 \\
21 & Particles I: case particles が, を, に, の & 10.1.1--10.1.4 \\
22 & Particles II: postpositions に, で, と, へ, から, まで; contrast between で and に & 10.2.1--10.2.6 \\
23 & Adverbial particles; も, しか, だけ, さえ, ばかり, ほど, くらい & 10.3 \\
24 & Sentence-final particles; か, ね, よ, よね and interactional stance & 10.4.1--10.4.4 \\
25 & Compound particles; として, について, に対して, によって, etc. & 10.5 \\
26 & は in greater detail; topic noun phrases, contrastive は, and は vs. が & 11.1--11.5 \\
27 & Noun types; common nouns, proper nouns, plural たち, and verbal nouns & 7.1.1; 7.1.2 \\
28 & Introductory formal nouns; こと, の, もの, ところ, わけ & 7.1.3; 18.1--18.6 \\
29 & Nominalization with こと and の; 読むこと and 読むの & 18.1; 18.2; 26.1 \\
30 & Explanatory のだ/んです & 18.2.2 \\
31 & Formal nouns and related expressions; もの, わけ, はず, つもり, ため, だけ, ばかり & 18.3; 18.6--18.9 \\
32 & Semantic classification of verbs; transitivity, action and state, punctual and durative verbs, controllability & 12.2--12.5 \\
33 & Potential forms and potential expressions; 読める, 食べられる, できる & 12.6 \\
34 & Giving and receiving verbs; あげる, くれる, もらう; viewpoint & 12.7 \\
35 & Clothing and wearing verbs; body and clothing-related verbal expressions & 12.8 \\
36 & ている; progressive, resultative state, experience, and habitual meaning & 19.2.1; 27.1--27.2 \\
37 & Auxiliary verbs with て; てある, てみる, ておく, てしまう & 19.2.2--19.2.5 \\
38 & てくる and ていく; direction, change, and temporal development & 19.2.6--19.2.7 \\
39 & て plus giving and receiving; てあげる, てくれる, てもらう & 19.2.8 \\
40 & Compound verbs based on the continuative form; 始める, 出す, 終わる, 過ぎる, 得る & 19.3.1--19.3.4 \\
41 & Causative forms; morphology, basic causative meaning, permission and coercion & 20.1; 20.2 \\
42 & Further causative patterns; adversative causative, causative with giving and receiving, lexical causatives, Aく/にする & 20.3--20.5 \\
43 & Passive forms; ordinary passive and によって passive & 21.1--21.3 \\
44 & Adversative passive, inanimate-subject passive, and causative-passive & 21.4; 21.5; 21.7 \\
45 & Clause connection I; coordination, sequence, and reason; て, し, たり, から, ので & 22.1--22.4 \\
46 & Clause connection II; contrast, concession, alternatives, and topic shift; けど, が, のに, ても, ながら, か, なり & 22.5--22.7 \\
47 & Temporal clauses; 時, 前, 後, 間, ところ; tense in subordinate clauses & 23.1--23.3; 27.4.2 \\
48 & Conditional clauses; と, ば, なら, たら, のなら & 24.1--24.4 \\
49 & Relative clauses and complement clauses; quotation with と/って; embedded questions & 25.1--25.2; 26.2--26.3 \\
50 & Evidentiality, conjecture, and honorific language; そう, よう, らしい; honorific, humble, polite, and beautifying forms & 28.1--28.5; 29.1--29.8 \\

\end{longtable}
}
\tableofcontents
\mainmatter
\chapter{1000語レベル}
\begin{readersection}{\ruby{学校}{がっこう}の\ruby{会}{かい}}
\voc[此処]{\ruby{此処}{ここ}}{here}は、\voc[県]{\ruby{県}{けん}}{prefecture}の\voc[学校]{\ruby{学校}{がっこう}}{school}です。\ruby{学校}{がっこう}の\voc[画面]{\ruby{画面}{がめん}}{screen}には、\voc[最初]{\ruby{最初}{さいしょ}}{first}の\voc[記事]{\ruby{記事}{きじ}}{article}があります。\ruby{記事}{きじ}は、\voc[高校]{\ruby{高校}{こうこう}}{high school}の\voc[会]{\ruby{会}{かい}}{meeting}の\ruby{案内}{あんない}です。この\ruby{会}{かい}は、\ruby{県}{けん}の\voc[事業]{\ruby{事業}{じぎょう}}{project}の\voc[一部]{\ruby{一部}{いちぶ}}{part}です。\ruby{会}{かい}の\voc[中心]{\ruby{中心}{ちゅうしん}}{center}は、\voc[安全]{\ruby{安全}{あんぜん}}{safety}と\voc[自然]{\ruby{自然}{しぜん}}{nature}です。\voc[目標]{\ruby{目標}{もくひょう}}{goal}は、\ruby{人々}{ひとびと}の\ruby{安全}{あんぜん}です。\ruby{安全}{あんぜん}は\voc[大事]{\ruby{大事}{だいじ}}{important}です。\ruby{自然}{しぜん}も\ruby{大事}{だいじ}です。

\ruby{学校}{がっこう}の\voc[壁]{\ruby{壁}{かべ}}{wall}には、大きい\voc[シート]{シート}{sheet}があります。シートの\voc[左]{\ruby{左}{ひだり}}{left}には、「\ruby{予定}{よてい}」の\ruby{欄}{らん}があります。\ruby{右}{みぎ}には、「\voc[予想]{\ruby{予想}{よそう}}{prediction}」の\ruby{欄}{らん}があります。\ruby{予定}{よてい}の\voc[期間]{\ruby{期間}{きかん}}{period}は一\voc[月]{\ruby{月}{かげつ}}{month}です。\ruby{期間}{きかん}の\voc[途中]{\ruby{途中}{とちゅう}}{middle}には、\voc[試合]{\ruby{試合}{しあい}}{match}があります。\ruby{試合}{しあい}は\voc[夜]{\ruby{夜}{よる}}{night}です。\ruby{夜}{よる}の\voc[移動]{\ruby{移動}{いどう}}{movement}は\ruby{安全}{あんぜん}です。\ruby{移動}{いどう}の\ruby{道}{みち}には\voc[電気]{\ruby{電気}{でんき}}{electricity}があります。\ruby{道}{みち}は\ruby{明}{あか}るいです。

\ruby{予想}{よそう}の\ruby{欄}{らん}には、\voc[パーセント]{パーセント}{percent}の\ruby{数字}{すうじ}があります。\voc[例えば]{\ruby{例えば}{たとえば}}{for example}、\ruby{安全}{あんぜん}の\ruby{数字}{すうじ}、\voc[人気]{\ruby{人気}{にんき}}{popularity}の\ruby{数字}{すうじ}、\voc[影響]{\ruby{影響}{えいきょう}}{influence}の\ruby{数字}{すうじ}です。\voc[発生]{\ruby{発生}{はっせい}}{occurrence}の\ruby{欄}{らん}もあります。この\ruby{発生}{はっせい}は、\voc[悪い]{\ruby{悪}{わる}い}{bad}ニュースではありません。\ruby{学校}{がっこう}の\ruby{小}{ちい}さい\ruby{問題}{もんだい}です。\ruby{先生}{せんせい}は、この\ruby{問題}{もんだい}を\voc[考える]{\ruby{考}{かんが}える}{to think}人です。

\ruby{会}{かい}の\ruby{人}{ひと}は、\voc[親]{\ruby{親}{おや}}{parent}、\voc[兄]{\ruby{兄}{あに}}{older brother}、\ruby{高校}{こうこう}の\ruby{学生}{がくせい}、そして\ruby{先生}{せんせい}です。\voc[貴方]{\ruby{貴方}{あなた}}{you}も、この\ruby{会}{かい}の\ruby{人}{ひと}です。\voc[誰]{\ruby{誰}{だれ}}{who}がこの\ruby{会}{かい}の\voc[主人]{\ruby{主人}{しゅじん}}{master}ですか。\voc[彼処]{\ruby{彼処}{あそこ}}{there}の\ruby{人}{ひと}ですか。いいえ、\ruby{主人}{しゅじん}ではありません。\voc[全員]{\ruby{全員}{ぜんいん}}{everyone}が、この\ruby{会}{かい}の\ruby{大事}{だいじ}な\ruby{人}{ひと}です。\voc[人々]{\ruby{人々}{ひとびと}}{people}は\voc[元気]{\ruby{元気}{げんき}}{well}です。\ruby{親}{おや}も\ruby{兄}{あに}も\ruby{学生}{がくせい}も、\ruby{共}{とも}に\ruby{会}{かい}の\ruby{人}{ひと}です。

\ruby{学校}{がっこう}の\ruby{入口}{いりぐち}には、\voc[ハブ]{ハブ}{hub}の\ruby{画面}{がめん}があります。このハブは、\ruby{会}{かい}の\ruby{情報}{じょうほう}の\ruby{中心}{ちゅうしん}です。\ruby{画面}{がめん}の\ruby{中}{なか}には、\voc[色んな]{\ruby{色}{いろ}んな}{various}\ruby{記事}{きじ}があります。\voc[島]{\ruby{島}{しま}}{island}の\ruby{自然}{しぜん}、\voc[未来]{\ruby{未来}{みらい}}{future}の\voc[価値]{\ruby{価値}{かち}}{value}、\voc[投資]{\ruby{投資}{とうし}}{investment}の\voc[展開]{\ruby{展開}{てんかい}}{development}、\voc[治療]{\ruby{治療}{ちりょう}}{treatment}の\voc[終了]{\ruby{終了}{しゅうりょう}}{end}、\voc[仕事]{\ruby{仕事}{しごと}}{work}の\ruby{記事}{きじ}、\voc[料理]{\ruby{料理}{りょうり}}{cooking}の\ruby{記事}{きじ}です。それぞれの\ruby{記事}{きじ}には、\voc[役]{\ruby{役}{やく}}{role}があります。\voc[神]{\ruby{神}{かみ}}{god}の\ruby{話}{はなし}もあります。それは、\voc[所謂]{\ruby{所謂}{いわゆる}}{so-called}\ruby{古}{ふる}い\ruby{話}{はなし}です。

\ruby{別}{べつ}の\ruby{欄}{らん}は、\ruby{物}{もの}の\ruby{欄}{らん}です。\ruby{欄}{らん}には、\voc[ドア]{ドア}{door}、\voc[カー]{カー}{car}、\voc[装備]{\ruby{装備}{そうび}}{equipment}、\voc[コーヒー]{コーヒー}{coffee}、\voc[肉]{\ruby{肉}{にく}}{meat}、そして\ruby{料理}{りょうり}の\ruby{写真}{しゃしん}があります。\voc[値段]{\ruby{値段}{ねだん}}{price}は\voc[安い]{\ruby{安}{やす}い}{cheap}\ruby{値段}{ねだん}です。\ruby{親}{おや}はコーヒーを\voc[買う]{\ruby{買}{か}う}{to buy}人です。\ruby{兄}{あに}は\ruby{肉}{にく}を\ruby{買}{か}う人です。\ruby{学生}{がくせい}はコーヒーを\voc[飲む]{\ruby{飲}{の}む}{to drink}人です。\ruby{肉}{にく}は\ruby{悪}{わる}い\ruby{肉}{にく}ですか。いいえ、\ruby{悪}{わる}い\ruby{肉}{にく}ではありません。\voc[痛い]{\ruby{痛}{いた}い}{painful}\ruby{話}{はなし}でもありません。\voc[怖い]{\ruby{怖}{こわ}い}{scary}\ruby{話}{はなし}でもありません。

\ruby{画面}{がめん}には、\voc[キャラ]{キャラ}{character}の\ruby{欄}{らん}もあります。このキャラは、\voc[姿]{\ruby{姿}{すがた}}{figure}が\ruby{面白}{おもしろ}いキャラです。\voc[そんな]{そんな}{such}キャラは、\voc[大分]{\ruby{大分}{だいぶ}}{quite}\ruby{人気}{にんき}です。でも、キャラの\ruby{人気}{にんき}は、この\ruby{会}{かい}の\ruby{中心}{ちゅうしん}ではありません。\voc[最も]{\ruby{最}{もっと}も}{most}\ruby{大事}{だいじ}な\ruby{欄}{らん}は、\ruby{安全}{あんぜん}の\ruby{欄}{らん}です。\voc[取り敢えず]{\ruby{取}{と}り\ruby{敢}{あ}えず}{for now}、この\ruby{画面}{がめん}は\voc[クリア]{クリア}{clear}です。どの\ruby{欄}{らん}も\ruby{見}{み}やすいです。

\ruby{先生}{せんせい}は、\ruby{会}{かい}の\ruby{仕事}{しごと}を\voc[進める]{\ruby{進}{すす}める}{to advance}人です。\ruby{学生}{がくせい}も、\ruby{仕事}{しごと}を\ruby{進}{すす}める人です。\ruby{親}{おや}は、\ruby{安全}{あんぜん}な\ruby{移動}{いどう}を\voc[勧める]{\ruby{勧}{すす}める}{to recommend}人です。\ruby{兄}{あに}は、\ruby{新}{あたら}しい\ruby{道}{みち}を\voc[探す]{\ruby{探}{さが}す}{to search}人です。\ruby{先生}{せんせい}は、\ruby{画面}{がめん}を\voc[見る]{\ruby{見}{み}る}{to see}人です。\ruby{誰}{だれ}かを\voc[待つ]{\ruby{待}{ま}つ}{to wait}人でもあります。\ruby{会}{かい}の\ruby{途中}{とちゅう}で\ruby{問題}{もんだい}が\voc[落ちる]{\ruby{落}{お}ちる}{to fall}こともあります。でも、\ruby{先生}{せんせい}は\voc[続ける]{\ruby{続}{つづ}ける}{to continue}人です。

この\ruby{学校}{がっこう}に\voc[住む]{\ruby{住}{す}む}{to live}人はいません。でも、\ruby{学校}{がっこう}で\ruby{長}{なが}く\voc[生きる]{\ruby{生}{い}きる}{to live}ような\ruby{思}{おも}い\ruby{出}{で}はあります。\ruby{学生}{がくせい}は、この\ruby{会}{かい}から\ruby{大事}{だいじ}な\ruby{経験}{けいけん}を\voc[得る]{\ruby{得}{え}る}{to obtain}人です。\ruby{古}{ふる}い\ruby{考}{かんが}えから\voc[抜ける]{\ruby{抜}{ぬ}ける}{to come out}人もいます。\ruby{先生}{せんせい}は、\ruby{会}{かい}の\ruby{形}{かたち}を\voc[変える]{\ruby{変}{か}える}{to change}人です。\ruby{新}{あたら}しい\ruby{意見}{いけん}を\voc[加える]{\ruby{加}{くわ}える}{to add}人でもあります。そして、\ruby{画面}{がめん}は、その\ruby{変化}{へんか}を\voc[示す]{\ruby{示}{しめ}す}{to show}ものです。

\ruby{会}{かい}には、\ruby{丁寧}{ていねい}な\ruby{言葉}{ことば}もあります。\ruby{先生}{せんせい}は、\voc[為さる]{\ruby{為}{な}さる}{to do}という\ruby{言葉}{ことば}を\ruby{使}{つか}います。\ruby{学生}{がくせい}は、\voc[頂く]{\ruby{頂}{いただ}く}{to receive}という\ruby{言葉}{ことば}を\ruby{見}{み}ます。\voc[御免]{\ruby{御免}{ごめん}}{sorry}という\ruby{言葉}{ことば}も、\ruby{会}{かい}の\ruby{途中}{とちゅう}で\ruby{出}{で}ます。\ruby{誰}{だれ}かが\voc[怒る]{\ruby{怒}{おこ}る}{to get angry}時、「\ruby{御免}{ごめん}」は\ruby{大事}{だいじ}な\ruby{言葉}{ことば}です。\voc[共]{\ruby{共}{とも}}{together}に\ruby{考}{かんが}えることも\ruby{大事}{だいじ}です。\voc[侭]{\ruby{侭}{まま}}{as is}にしないで、\ruby{問題}{もんだい}を\ruby{変}{か}えることも\ruby{大事}{だいじ}です。

\ruby{先生}{せんせい}の\ruby{最後}{さいご}のメモには、\voc[完成]{\ruby{完成}{かんせい}}{completion}と\ruby{終了}{しゅうりょう}の\ruby{言葉}{ことば}があります。\ruby{会}{かい}の\ruby{準備}{じゅんび}は、まだ\ruby{完成}{かんせい}ではありません。\ruby{事業}{じぎょう}の\ruby{終了}{しゅうりょう}もまだです。\voc[早速]{\ruby{早速}{さっそく}}{at once}、\voc[クリック]{クリック}{click}です。クリックの\voc[先]{\ruby{先}{さき}}{ahead}には、\ruby{次}{つぎ}の\ruby{記事}{きじ}があります。その\ruby{記事}{きじ}は、\ruby{未来}{みらい}の\ruby{会}{かい}の\ruby{記事}{きじ}です。\ruby{未来}{みらい}の\ruby{価値}{かち}は、まだ\ruby{予想}{よそう}だけです。でも、\ruby{人々}{ひとびと}は\ruby{共}{とも}に\ruby{考}{かんが}えます。

この\ruby{第一課}{だいいっか}は\ruby{長}{なが}いです。でも、\ruby{第一課}{だいいっか}の\ruby{中心}{ちゅうしん}は、\ruby{名詞}{めいし}の\ruby{文}{ぶん}です。\ruby{此処}{ここ}は\ruby{学校}{がっこう}です。\ruby{画面}{がめん}は\ruby{学校}{がっこう}の\ruby{画面}{がめん}です。\ruby{記事}{きじ}は\ruby{高校}{こうこう}の\ruby{会}{かい}の\ruby{記事}{きじ}です。\ruby{安全}{あんぜん}は\ruby{大事}{だいじ}です。\ruby{自然}{しぜん}も\ruby{大事}{だいじ}です。これが\ruby{第一課}{だいいっか}の\ruby{基礎}{きそ}です。
\end{readersection}

\begin{readersection}{\ruby{街}{まち}の\ruby{小}{ちい}さな\ruby{会}{かい}}
\voc[今回]{\ruby{今回}{こんかい}}{this time}の\ruby{会}{かい}は、\ruby{学校}{がっこう}の\ruby{会}{かい}ではありません。\voc[街]{\ruby{街}{まち}}{town}の\ruby{会}{かい}です。\ruby{場所}{ばしょ}は、\ruby{学校}{がっこう}の\ruby{前}{まえ}の\voc[ホテル]{ホテル}{hotel}です。\voc[前]{\ruby{前}{まえ}}{front}には、\ruby{大}{おお}きい\ruby{画面}{がめん}と\ruby{白}{しろ}いシートがあります。シートの\ruby{真}{ま}ん\ruby{中}{なか}には、\voc[今回]{\ruby{今回}{こんかい}}{this time}の\voc[内容]{\ruby{内容}{ないよう}}{content}があります。\ruby{内容}{ないよう}は\voc[小さい]{\ruby{小}{ちい}さい}{small}ですが、\voc[意味]{\ruby{意味}{いみ}}{meaning}は大きいです。

\ruby{画面}{がめん}の\voc[真ん中]{\ruby{真}{ま}ん\ruby{中}{なか}}{middle}には、\voc[小さな]{\ruby{小}{ちい}さな}{small}\ruby{記事}{きじ}があります。その\ruby{記事}{きじ}は、\ruby{街}{まち}の\voc[歴史]{\ruby{歴史}{れきし}}{history}の\ruby{記事}{きじ}です。\voc[昔]{\ruby{昔}{むかし}}{past}、この\ruby{街}{まち}は\ruby{静}{しず}かな\ruby{場所}{ばしょ}でした。今も、\ruby{山}{やま}の\ruby{前}{まえ}の\ruby{場所}{ばしょ}は\ruby{静}{しず}かです。でも、\voc[最近]{\ruby{最近}{さいきん}}{recently}は、\ruby{街}{まち}の\ruby{中心}{ちゅうしん}が\ruby{大分}{だいぶ}にぎやかです。\ruby{学校}{がっこう}の\ruby{学生}{がくせい}も、\ruby{親}{おや}も、\ruby{兄}{あに}も、ここに\ruby{来}{き}ます。

\ruby{会}{かい}の\ruby{目的}{もくてき}は\ruby{安全}{あんぜん}です。\voc[目的]{\ruby{目的}{もくてき}}{purpose}は\ruby{大事}{だいじ}です。\ruby{夜}{よる}の\ruby{移動}{いどう}は、\voc[大丈夫]{\ruby{大丈夫}{だいじょうぶ}}{okay}ですか。\voc[正直]{\ruby{正直}{しょうじき}}{honestly}、まだ\ruby{少}{すこ}し\voc[弱い]{\ruby{弱}{よわ}い}{weak}です。\ruby{電気}{でんき}はありますが、\ruby{古}{ふる}いドアの\ruby{前}{まえ}は\ruby{暗}{くら}いです。\ruby{先生}{せんせい}は「この\ruby{場所}{ばしょ}は\voc[嫌]{\ruby{嫌}{いや}}{unpleasant}ではありません。でも、\ruby{少}{すこ}し\ruby{危}{あぶ}ないです」と\ruby{言}{い}います。\ruby{親}{おや}は、その\ruby{意見}{いけん}に\voc[反対]{\ruby{反対}{はんたい}}{opposition}ではありません。

\ruby{画面}{がめん}の\ruby{右}{みぎ}の\voc[欄]{\ruby{欄}{らん}}{column}には、\ruby{会}{かい}の\ruby{予定}{よてい}があります。\ruby{最初}{さいしょ}は\ruby{安全}{あんぜん}の\ruby{話}{はなし}です。\ruby{次}{つぎ}は\ruby{街}{まち}の\ruby{歴史}{れきし}です。\voc[そして]{そして}{and then}、\ruby{最後}{さいご}は\ruby{質問}{しつもん}です。\voc[質問]{\ruby{質問}{しつもん}}{question}の\ruby{時間}{じかん}は一\voc[週間]{\ruby{週間}{しゅうかん}}{week}ではありません。一\ruby{時間}{じかん}です。\ruby{一}{いち}\voc[秒]{\ruby{秒}{びょう}}{second}だけの\ruby{質問}{しつもん}でも\ruby{大丈夫}{だいじょうぶ}です。\voc[一]{\ruby{一}{いち}}{one}つの\ruby{質問}{しつもん}も、\ruby{大事}{だいじ}な\ruby{質問}{しつもん}です。

\ruby{会}{かい}の\ruby{前}{まえ}に、\ruby{先生}{せんせい}は\ruby{画面}{がめん}を\voc[直す]{\ruby{直}{なお}す}{to fix}人です。\ruby{学生}{がくせい}は\ruby{机}{つくえ}の\ruby{上}{うえ}に\voc[パソコン]{パソコン}{computer}を\ruby{置}{お}きます。\voc[カメラ]{カメラ}{camera}もあります。カメラの\voc[画像]{\ruby{画像}{がぞう}}{image}は\ruby{画面}{がめん}にあります。この\ruby{画像}{がぞう}は\ruby{小}{ちい}さいですが、\ruby{色}{いろ}が\ruby{明}{あか}るいです。\voc[色]{\ruby{色}{いろ}}{color}は\ruby{大事}{だいじ}です。\ruby{古}{ふる}い\ruby{画面}{がめん}の\ruby{色}{いろ}は\ruby{弱}{よわ}いです。\ruby{新}{あたら}しい\ruby{画面}{がめん}の\ruby{色}{いろ}は\ruby{強}{つよ}いです。

\ruby{先生}{せんせい}は、\ruby{会}{かい}の\ruby{資料}{しりょう}を\voc[作成]{\ruby{作成}{さくせい}}{creation}する人です。\ruby{資料}{しりょう}の\voc[セット]{セット}{set}は三つです。一つは\ruby{歴史}{れきし}のセットです。一つは\ruby{安全}{あんぜん}のセットです。一つは\ruby{街}{まち}の\ruby{写真}{しゃしん}のセットです。\ruby{写真}{しゃしん}の\voc[バック]{バック}{back}は\ruby{山}{やま}です。\voc[山]{\ruby{山}{やま}}{mountain}は\ruby{大}{おお}きいです。\ruby{山}{やま}の\ruby{色}{いろ}は\ruby{古}{ふる}い\ruby{写真}{しゃしん}では\ruby{暗}{くら}いですが、\ruby{今回}{こんかい}の\ruby{画像}{がぞう}では\ruby{明}{あか}るいです。

\ruby{会}{かい}には、\ruby{街}{まち}の\ruby{人々}{ひとびと}も\ruby{来}{き}ます。\voc[客]{\ruby{客}{きゃく}}{guest}も\ruby{来}{き}ます。\voc[彼]{\ruby{彼}{かれ}}{he}は\ruby{昔}{むかし}この\ruby{街}{まち}に\voc[住む]{\ruby{住}{す}む}{to live}人でした。\ruby{彼}{かれ}は\ruby{今}{いま}、\ruby{別}{べつ}の\ruby{地域}{ちいき}にいます。\voc[地域]{\ruby{地域}{ちいき}}{region}は\ruby{遠}{とお}いです。でも、\ruby{彼}{かれ}は\ruby{明日}{あした}ここに\ruby{来}{き}ます。\voc[明日]{\ruby{明日}{あした}}{tomorrow}の\ruby{天気}{てんき}は\ruby{大丈夫}{だいじょうぶ}\voc[そう]{そう}{seem}です。\voc[成る程]{\ruby{成}{な}る\ruby{程}{ほど}}{I see}、これは\ruby{嬉}{うれ}しいニュースです。\voc[嬉しい]{\ruby{嬉}{うれ}しい}{happy}ニュースです。

\ruby{彼}{かれ}の\ruby{話}{はなし}は\voc[面白い]{\ruby{面白}{おもしろ}い}{interesting}です。\voc[実]{\ruby{実}{じつ}}{truth}は、\ruby{彼}{かれ}はこの\ruby{街}{まち}で\voc[生まれる]{\ruby{生}{う}まれる}{to be born}人です。\ruby{昔}{むかし}の\ruby{学校}{がっこう}、\ruby{古}{ふる}い\ruby{壁}{かべ}、\ruby{小}{ちい}さいドア、\ruby{山}{やま}の\ruby{前}{まえ}の\ruby{道}{みち}、それらは\ruby{彼}{かれ}の\voc[経験]{\ruby{経験}{けいけん}}{experience}です。\ruby{彼}{かれ}は、その\ruby{経験}{けいけん}を\voc[語る]{\ruby{語}{かた}る}{to tell}人です。\ruby{学生}{がくせい}は、その\ruby{話}{はなし}に\voc[興味]{\ruby{興味}{きょうみ}}{interest}があります。\ruby{親}{おや}も\ruby{兄}{あに}も\ruby{興味}{きょうみ}があります。

\ruby{会}{かい}の\ruby{途中}{とちゅう}で、\ruby{学生}{がくせい}は\ruby{先生}{せんせい}に\ruby{質問}{しつもん}します。「\voc[一体]{\ruby{一体}{いったい}}{on earth}、この\ruby{会}{かい}の\ruby{本当}{ほんとう}の\ruby{意味}{いみ}は\ruby{何}{なん}ですか。」\ruby{先生}{せんせい}は\ruby{少}{すこ}し\ruby{待}{ま}ちます。\voc[直ぐ]{\ruby{直}{す}ぐ}{soon}には\ruby{答}{こた}えません。\voc[兎に角]{\ruby{兎}{と}に\ruby{角}{かく}}{anyway}、\ruby{質問}{しつもん}は\ruby{大事}{だいじ}です。\ruby{先生}{せんせい}は\ruby{笑}{わら}って、「この\ruby{会}{かい}の\ruby{意味}{いみ}は、\ruby{街}{まち}を\ruby{知}{し}ることです」と\ruby{言}{い}います。

\ruby{学生}{がくせい}は\ruby{更}{さら}に\ruby{質問}{しつもん}します。「\ruby{古}{ふる}い\ruby{歴史}{れきし}だけですか。\ruby{未来}{みらい}は\ruby{大事}{だいじ}ではありませんか。」\ruby{先生}{せんせい}は「\ruby{未来}{みらい}も\ruby{大事}{だいじ}です」と\ruby{言}{い}います。\ruby{街}{まち}の\ruby{未来}{みらい}には、\voc[利用]{\ruby{利用}{りよう}}{use}の\ruby{問題}{もんだい}があります。\ruby{画面}{がめん}の\ruby{利用}{りよう}、\ruby{学校}{がっこう}の\ruby{利用}{りよう}、\ruby{道}{みち}の\ruby{利用}{りよう}です。\ruby{政府}{せいふ}の\ruby{記事}{きじ}もあります。\voc[政府]{\ruby{政府}{せいふ}}{government}の\ruby{対応}{たいおう}は\ruby{早}{はや}いですか。\voc[対応]{\ruby{対応}{たいおう}}{response}はまだ\ruby{十分}{じゅうぶん}ではありません。

\ruby{会}{かい}の\ruby{後}{あと}で、\ruby{学生}{がくせい}は\ruby{街}{まち}を\ruby{歩}{ある}きます。ホテルの\ruby{前}{まえ}に\voc[犬]{\ruby{犬}{いぬ}}{dog}が\voc[居る]{いる}{to be}。その\ruby{犬}{いぬ}は\ruby{小}{ちい}さいです。でも、\ruby{声}{こえ}は\ruby{大}{おお}きいです。犬の\ruby{前}{まえ}には\ruby{赤}{あか}い\ruby{服}{ふく}の\ruby{人}{ひと}がいます。\voc[服]{\ruby{服}{ふく}}{clothes}は\ruby{新}{あたら}しいです。\ruby{色}{いろ}も\ruby{明}{あか}るいです。ホテルの\ruby{客}{きゃく}は、その\ruby{服}{ふく}を\ruby{見}{み}ます。\ruby{学生}{がくせい}も\ruby{見}{み}ます。\ruby{先生}{せんせい}は「その\ruby{色}{いろ}は\ruby{会}{かい}の\ruby{画面}{がめん}に\ruby{合}{あ}います」と\ruby{言}{い}います。

\ruby{次}{つぎ}の日、\ruby{学生}{がくせい}は\ruby{資料}{しりょう}を\voc[公開]{\ruby{公開}{こうかい}}{release}します。\ruby{公開}{こうかい}の\ruby{前}{まえ}に、\ruby{先生}{せんせい}は\ruby{画像}{がぞう}を\voc[外す]{\ruby{外}{はず}す}{to remove}人です。\ruby{古}{ふる}い\ruby{写真}{しゃしん}を\ruby{外}{はず}します。そして、\ruby{新}{あたら}しい\ruby{画像}{がぞう}を\voc[付ける]{\ruby{付}{つ}ける}{to attach}人です。\ruby{学生}{がくせい}は、\ruby{公開}{こうかい}の\ruby{日}{ひ}を\voc[スタート]{スタート}{start}の日にします。\ruby{会}{かい}の\ruby{準備}{じゅんび}は\voc[既に]{\ruby{既}{すで}に}{already}\ruby{完成}{かんせい}です。\voc[一旦]{\ruby{一旦}{いったん}}{once}、\ruby{画面}{がめん}を\ruby{確認}{かくにん}します。

\ruby{公開}{こうかい}の\ruby{画面}{がめん}には、\ruby{新}{あたら}しい\voc[モード]{モード}{mode}があります。\ruby{安全}{あんぜん}モード、\ruby{歴史}{れきし}モード、\ruby{質問}{しつもん}モードです。\ruby{学生}{がくせい}は\ruby{安全}{あんぜん}モードを\voc[選択]{\ruby{選択}{せんたく}}{choice}します。\ruby{歴史}{れきし}モードも\ruby{面白}{おもしろ}いですが、\ruby{安全}{あんぜん}モードは\ruby{今回}{こんかい}の\ruby{中心}{ちゅうしん}です。\ruby{画面}{がめん}の\ruby{レベル}{レベル}は\voc[レベル]{レベル}{level}一です。\ruby{第一課}{だいいっか}より\ruby{少}{すこ}し\ruby{高}{たか}いです。でも、まだ\ruby{無理}{むり}ではありません。\voc[無理]{\ruby{無理}{むり}}{impossible}ではない\ruby{内容}{ないよう}です。

\ruby{昼}{ひる}、\ruby{学生}{がくせい}は\ruby{街}{まち}の\ruby{小}{ちい}さい\ruby{店}{みせ}に\voc[乗る]{\ruby{乗}{の}る}{to ride}バスで\ruby{行}{い}きます。バスの\ruby{値段}{ねだん}は五百\voc[円]{\ruby{円}{えん}}{yen}です。\ruby{円}{えん}は\ruby{安}{やす}いですか。\ruby{学生}{がくせい}には\ruby{少}{すこ}し\ruby{高}{たか}いです。店の\ruby{料理}{りょうり}は\voc[美味しい]{\ruby{美味}{おい}しい}{delicious}です。コーヒーも\ruby{美味}{おい}しいです。でも、\ruby{先生}{せんせい}は「\ruby{値段}{ねだん}が\ruby{少}{すこ}し\voc[上がる]{\ruby{上}{あ}がる}{to rise}かもしれません」と\ruby{言}{い}います。\ruby{学生}{がくせい}は「\ruby{値段}{ねだん}が\ruby{下}{さ}がると\ruby{嬉}{うれ}しいです」と\ruby{言}{い}います。\voc[下がる]{\ruby{下}{さ}がる}{to fall}ことは\ruby{嬉}{うれ}しいことです。

\ruby{会}{かい}の\ruby{終}{お}わりに、\ruby{先生}{せんせい}は\ruby{学生}{がくせい}に\ruby{電話}{でんわ}をします。\voc[電話]{\ruby{電話}{でんわ}}{phone}の\ruby{声}{こえ}は\ruby{少}{すこ}し\ruby{遅}{おそ}いです。\voc[遅い]{\ruby{遅}{おそ}い}{slow}ですが、\ruby{内容}{ないよう}は\ruby{分}{わ}かります。\ruby{先生}{せんせい}は「\ruby{今回}{こんかい}の\ruby{会}{かい}は\ruby{中々}{なかなか}よいです」と\ruby{言}{い}います。\voc[中々]{\ruby{中々}{なかなか}}{quite}\ruby{良}{よ}い\ruby{会}{かい}です。\voc[詰まり]{\ruby{詰}{つ}まり}{in short}、\ruby{会}{かい}は\ruby{成功}{せいこう}です。\ruby{学生}{がくせい}は\ruby{嬉}{うれ}しいです。\ruby{親}{おや}も\ruby{嬉}{うれ}しいです。\ruby{兄}{あに}も\ruby{嬉}{うれ}しいです。

\ruby{最後}{さいご}に、\ruby{学生}{がくせい}は\ruby{短}{みじか}い\ruby{文}{ぶん}を\ruby{書}{か}きます。「\ruby{今回}{こんかい}の\ruby{会}{かい}は\ruby{街}{まち}の\ruby{会}{かい}です。\ruby{場所}{ばしょ}はホテルです。\ruby{内容}{ないよう}は\ruby{安全}{あんぜん}と\ruby{歴史}{れきし}です。\ruby{質問}{しつもん}は\ruby{多}{おお}いです。でも、\ruby{対応}{たいおう}は\ruby{大丈夫}{だいじょうぶ}です。\ruby{人々}{ひとびと}は\ruby{元気}{げんき}です。\ruby{未来}{みらい}の\ruby{価値}{かち}もあります。\ruby{何}{いず}れ、この\ruby{街}{まち}はもっと\ruby{良}{よ}い\ruby{街}{まち}です。」\voc[何れ]{\ruby{何}{いず}れ}{eventually}、この\ruby{会}{かい}の\ruby{意味}{いみ}はもっと\ruby{知}{し}れます。\voc[知れる]{\ruby{知}{し}れる}{to be known}ことは、\ruby{街}{まち}にとって\ruby{幸}{しあわ}せです。\voc[幸せ]{\ruby{幸}{しあわ}せ}{happiness}な\ruby{会}{かい}です。\voc[イエス]{イエス}{Jesus}の\ruby{話}{はなし}は、この\ruby{会}{かい}の\ruby{中心}{ちゅうしん}ではありません。でも、\ruby{神}{かみ}の\ruby{話}{はなし}の\ruby{中}{なか}に、イエスの\ruby{名前}{なまえ}も\ruby{少}{すこ}しあります。
\end{readersection}
\end{document}
